\documentclass[10pt,a4paper]{amsart}
\usepackage{amsmath,amssymb,amsthm,graphicx}
\usepackage{verbatim}
\newenvironment{reference}{\comment}{\endcomment}
%\newenvironment{reference}{}{}
\newenvironment{slogan}{\comment}{\endcomment}
\newenvironment{history}{\comment}{\endcomment}
% For commutative diagrams we use Xy-pic
\usepackage[all]{xy}

% We use 2cell for 2-commutative diagrams.
\xyoption{2cell}
% Theorem environments.
%
\theoremstyle{plain}
\newtheorem{theorem}[subsection]{Theorem}
\newtheorem{proposition}[subsection]{Proposition}
\newtheorem{lemma}[subsection]{Lemma}

\theoremstyle{definition}
\newtheorem{definition}[subsection]{Definition}
\newtheorem{example}[subsection]{Example}
\newtheorem{exercise}[subsection]{Exercise}
\newtheorem{situation}[subsection]{Situation}

\theoremstyle{remark}
\newtheorem{remark}[subsection]{Remark}
\newtheorem{remarks}[subsection]{Remarks}

\numberwithin{equation}{subsection}

% Macros
%
\def\lim{\mathop{\mathrm{lim}}\nolimits}
\def\colim{\mathop{\mathrm{colim}}\nolimits}
\def\Spec{\mathop{\mathrm{Spec}}}
\def\Hom{\mathop{\mathrm{Hom}}\nolimits}
\def\Ext{\mathop{\mathrm{Ext}}\nolimits}
\def\SheafHom{\mathop{\mathcal{H}\!\mathit{om}}\nolimits}
\def\SheafExt{\mathop{\mathcal{E}\!\mathit{xt}}\nolimits}
\def\Sch{\mathit{Sch}}
\def\Mor{\mathop{\mathrm{Mor}}\nolimits}
\def\Ob{\mathop{\mathrm{Ob}}\nolimits}
\def\Sh{\mathop{\mathit{Sh}}\nolimits}
\def\NL{\mathop{N\!L}\nolimits}
\def\CH{\mathop{\mathrm{CH}}\nolimits}
\def\proetale{{pro\text{-}\acute{e}tale}}
\def\etale{{\acute{e}tale}}
\def\QCoh{\mathit{QCoh}}
\def\Ker{\mathop{\mathrm{Ker}}}
\def\Im{\mathop{\mathrm{Im}}}
\def\Coker{\mathop{\mathrm{Coker}}}
\def\Coim{\mathop{\mathrm{Coim}}}

% Boxtimes
%
\DeclareMathSymbol{\boxtimes}{\mathbin}{AMSa}{"02}

%
% Macros for moduli stacks/spaces
%
\def\QCohstack{\mathcal{QC}\!\mathit{oh}}
\def\Cohstack{\mathcal{C}\!\mathit{oh}}
\def\Spacesstack{\mathcal{S}\!\mathit{paces}}
\def\Quotfunctor{\mathrm{Quot}}
\def\Hilbfunctor{\mathrm{Hilb}}
\def\Curvesstack{\mathcal{C}\!\mathit{urves}}
\def\Polarizedstack{\mathcal{P}\!\mathit{olarized}}
\def\Complexesstack{\mathcal{C}\!\mathit{omplexes}}
% \Pic is the operator that assigns to X its picard group, usage \Pic(X)
% \Picardstack_{X/B} denotes the Picard stack of X over B
% \Picardfunctor_{X/B} denotes the Picard functor of X over B
\def\Pic{\mathop{\mathrm{Pic}}\nolimits}
\def\Picardstack{\mathcal{P}\!\mathit{ic}}
\def\Picardfunctor{\mathrm{Pic}}
\def\Deformationcategory{\mathcal{D}\!\mathit{ef}}
\usepackage{cleveref}
\crefname{originallemma}{original lemma}{original lemmas}
\Crefname{originallemma}{Original lemma}{Original lemmas}
\makeatletter
\newlabel{lemma-additivity-divisors-restricted}{{02QG}{source}{}{sourcepointer-1}{}}
\newlabel{lemma-additivity-divisors-restricted@cref}{{[originallemma][0][]02QG}{[1][0][]source}{}{}{}}
\makeatother
\crefname{originalreference}{original reference}{original references}
\Crefname{originalreference}{Original reference}{Original references}
\makeatletter
\newlabel{situation-setup}{{02QL}{source}{}{sourcepointer-2}{}}
\newlabel{situation-setup@cref}{{[originalreference][0][]02QL}{[1][0][]source}{}{}{}}
\makeatother
\title{The cycle of the zero scheme of a section is the sum over the components of the cycles of the restricted sections}
\author{The Stacks Project Authors}
\date{stacks-project:chow.tex:local-snapshot-20261012; GNU FDL 1.2 or later}
\begin{document}
\maketitle
\noindent\textit{Editorial note.} All delivered mathematical text is the author's own, copied from the pinned Stacks Project chapter file chow.tex; no statement or proof step was rewritten, reordered or removed.\par\medskip
\noindent\textit{Editorial note.} Preamble provenance: the pinned chapter file opens with \textbackslash{}input\{preamble\} and carries no preamble text of its own, so the project's single Stacks-wide preamble staged at sources/stacks-chow/preamble.tex is excerpted verbatim as three spans, lines 8-12 (the verbatim package together with the project's own reference, slogan and history display-environment definitions, including the author's commented-out older form of the reference environment), lines 14-18 (Xy-pic setup) and lines 154-217 (the native theorem declarations and mathematical macros). The lines 8-12 span is required because a counted statement or its surrounding text may use those display environments, which the project typesets as comments. Nothing was added, removed or redefined; the wrapper class and the amsmath/amssymb/amsthm/graphicx base packages are supplied by the assembler.\par\medskip
\noindent\textit{Editorial note.} Same-chapter pointers are located at their \textbackslash{}label lines inside the pinned chapter file; cross-chapter pointers are located in the project's own tag list staged in the snapshot (sources/stacks-chow/tags.txt), at the verbatim lines carrying their canonical tags.\par\medskip
\noindent\textit{Allowed standard prerequisites (editorial).}\begin{itemize}
\item Additivity of divisors for a restricted family. Same chapter chow.tex, label lemma-additivity-divisors-restricted, invoked at proof line 4419. Supplies the final identity expressing a length over the local ring as the sum over its minimal primes of the local length times the length after quotienting by the prime, which is exactly the coefficient comparison the proof needs.
\end{itemize}\medskip
\ifdefined\hypertarget\hypertarget{sourcepointer-1}{}\fi
\noindent\textit{Source reference.} Original lemma 02QG: sources/stacks-chow/pinned-original-source.tex, source lines 373--373. Navigation pointer for \textbackslash{}ref\{lemma-additivity-divisors-restricted\} used inside the counted span (line 4420); the argument that consumes it is retained in the span.\par\medskip
\ifdefined\hypertarget\hypertarget{sourcepointer-2}{}\fi
\noindent\textit{Source reference.} Original reference 02QL: sources/stacks-chow/pinned-original-source.tex, source lines 1372--1372. Navigation pointer for \textbackslash{}ref\{situation-setup\} used inside the counted span (line 4359); the argument that consumes it is retained in the span.\par\medskip
\begin{lemma}
\label{lemma-prepare-geometric-cap}
Let $(S, \delta)$ be as in Situation \ref{situation-setup}.
Let $Y$ be locally of finite type over $S$.
Let $\mathcal{L}$ be an invertible $\mathcal{O}_Y$-module.
Let $s \in \Gamma(Y, \mathcal{L})$.
Assume
\begin{enumerate}
\item $\dim_\delta(Y) \leq k + 1$,
\item $\dim_\delta(Z(s)) \leq k$, and
\item for every generic point $\xi$ of an irreducible component of
$Z(s)$ of $\delta$-dimension $k$ the multiplication by $s$
induces an injection $\mathcal{O}_{Y, \xi} \to \mathcal{L}_\xi$.
\end{enumerate}
Write $[Y]_{k + 1} = \sum n_i[Y_i]$ where $Y_i \subset Y$ are the
irreducible components of $Y$ of $\delta$-dimension $k + 1$.
Set $s_i = s|_{Y_i} \in \Gamma(Y_i, \mathcal{L}|_{Y_i})$. Then
\begin{equation}
\label{equation-equal-as-cycles}
[Z(s)]_k =  \sum n_i[Z(s_i)]_k
\end{equation}
as $k$-cycles on $Y$.
\end{lemma}

\begin{proof}
Let $Z \subset Y$ be an integral closed subscheme of
$\delta$-dimension $k$. Let $\xi \in Z$ be its generic point.
We want to compare the coefficient $n$ of $[Z]$ in the expression
$\sum n_i[Z(s_i)]_k$ with the coefficient $m$ of $[Z]$ in the
expression $[Z(s)]_k$. Choose a generator $s_\xi \in \mathcal{L}_\xi$.
Write $A = \mathcal{O}_{Y, \xi}$, $L = \mathcal{L}_\xi$.
Then $L = As_\xi$. Write $s = f s_\xi$ for some (unique) $f \in A$.
Hypothesis (3) means that $f : A \to A$ is injective.
Since $\dim_\delta(Y) \leq k + 1$ and $\dim_\delta(Z) = k$
we have $\dim(A) = 0$ or $1$. We have
$$
m = \text{length}_A(A/(f))
$$
which is finite in either case.

\medskip\noindent
If $\dim(A) = 0$, then $f : A \to A$ being injective
implies that $f \in A^*$. Hence in this case $m$ is zero.
Moreover, the condition $\dim(A) = 0$ means that $\xi$
does not lie on any irreducible component of $\delta$-dimension
$k + 1$, i.e., $n = 0$ as well.

\medskip\noindent
Now, let $\dim(A) = 1$.
Since $A$ is a Noetherian local ring it has finitely
many minimal primes $\mathfrak q_1, \ldots, \mathfrak q_t$.
These correspond 1-1 with the $Y_i$ passing through $\xi'$.
Moreover $n_i = \text{length}_{A_{\mathfrak q_i}}(A_{\mathfrak q_i})$.
Also, the multiplicity of $[Z]$ in $[Z(s_i)]_k$ is
$\text{length}_A(A/(f, \mathfrak q_i))$.
Hence the equation to prove in this case is
$$
\text{length}_A(A/(f))
=
\sum \text{length}_{A_{\mathfrak q_i}}(A_{\mathfrak q_i})
\text{length}_A(A/(f, \mathfrak q_i))
$$
which follows from
Lemma \ref{lemma-additivity-divisors-restricted}.
\end{proof}

\end{document}
